% Generado por sistema/backend ReporteTex.java -- no editar a mano.
% Datos: /private/tmp/claude-501/-Users-miguel-project-taller-s/a868368b-4351-4d87-8017-fb9a5bc093aa/scratchpad/titulados.csv (titulados) / /private/tmp/claude-501/-Users-miguel-project-taller-s/a868368b-4351-4d87-8017-fb9a5bc093aa/scratchpad/empleadores.csv (empleadores)
\subsubsection*{C9. Experiencia profesional $\times$ remuneración}

\textit{Paso 1. Describir el problema.} Establecer si la remuneración depende de la experiencia del titulado (junior o consolidado, con corte en la mediana de los años de vida profesional), para conocer el grado de madurez profesional de la muestra.

\begin{itemize}[nosep]
  \item Variable independiente $X$: Experiencia.
  \item Variable dependiente $Y$: Remuneración.
  \item Población: trabajan en una organización ($N = 106$). Criterio ARCU-SUR: 3.2.3. Nivel de significancia $\alpha = 0.05$.
  \item Junior: hasta 6 años de vida profesional (la mediana); consolidado: más de 6 años.
  \item $H_0$: Las variables «Experiencia» (X) y «Remuneración» (Y) son independientes.
  \item $H_1$: Existe dependencia significativa entre las variables «Experiencia» (X) y «Remuneración» (Y).
\end{itemize}

\textit{Datos.} Primeros 10 de los 106 titulados que respondieron ambas preguntas (sin datos personales).

\begin{mitabla}{C9: datos de los primeros titulados}
\begin{tabular}{lll}
\toprule
Titulado & $X$: Experiencia & $Y$: Remuneración \\
\midrule
T1 & Junior & Más de Bs 5.000 \\
T2 & Consolidado & Más de Bs 5.000 \\
T3 & Junior & Hasta Bs 5.000 \\
T4 & Consolidado & Más de Bs 5.000 \\
T5 & Junior & Hasta Bs 5.000 \\
T6 & Consolidado & Más de Bs 5.000 \\
T7 & Junior & Más de Bs 5.000 \\
T8 & Consolidado & Más de Bs 5.000 \\
T9 & Consolidado & Más de Bs 5.000 \\
T10 & Consolidado & Más de Bs 5.000 \\
\bottomrule
\end{tabular}
\end{mitabla}

\textit{Paso 2. Tabla de contingencia} (frecuencia observada $FO_{ij}$, $N = 106$). Cada celda cuenta a los titulados que respondieron a la vez la fila ($X$) y la columna ($Y$).

\begin{mitabla}{C9: tabla de contingencia, frecuencia observada $FO_{ij}$}
\begin{tabular}{L{3.8cm}C{3.00cm}C{3.00cm}|c}
\toprule
$X$: Experiencia \textbackslash{} $Y$: Remuneración & Hasta Bs 5.000 & Más de Bs 5.000 & Total \\
\midrule
Junior & 30 & 23 & 53 \\
Consolidado & 3 & 50 & 53 \\
\midrule
Total & 33 & 73 & $N = 106$ \\
\bottomrule
\end{tabular}
\end{mitabla}

\textbf{Mapa de calor.} La misma tabla de contingencia, pintada por color en vez de por número:

\begin{mitabla}{C9: mapa de calor de $FO_{ij}$}
\begin{tabular}{L{3.8cm}C{3.00cm}C{3.00cm}}
\toprule
$X$: Experiencia \textbackslash{} $Y$: Remuneración & Hasta Bs 5.000 & Más de Bs 5.000 \\
\midrule
Junior & \cellcolor{gray!45}30 & \cellcolor{gray!35}23 \\
Consolidado & \cellcolor{gray!5}3 & \cellcolor{gray!75}50 \\
\bottomrule
\end{tabular}
\end{mitabla}

\textit{Paso 3. Frecuencia esperada.} Es lo que habría en cada celda si $X$ e $Y$ fueran independientes. Por ejemplo, para la celda Junior y Hasta Bs 5.000:
\[
FE_{11} = \frac{(53)(33)}{106} = 16.50
\]

\begin{mitabla}{C9: frecuencia esperada $FE_{ij}$}
\begin{tabular}{L{3.8cm}C{3.00cm}C{3.00cm}}
\toprule
$X$: Experiencia \textbackslash{} $Y$: Remuneración & Hasta Bs 5.000 & Más de Bs 5.000 \\
\midrule
Junior & $53\cdot33/106 = 16.50$ & $53\cdot73/106 = 36.50$ \\
Consolidado & $53\cdot33/106 = 16.50$ & $53\cdot73/106 = 36.50$ \\
\bottomrule
\end{tabular}
\end{mitabla}

\textit{Paso 4. Estadístico chi-cuadrado.} En cada celda se calcula $\frac{(FO_{ij}-FE_{ij})^2}{FE_{ij}}$ y luego se suman todas.

\begin{mitabla}{C9: cálculo de $X_0^2$ celda por celda}
\begin{tabular}{L{3.2cm}C{3.20cm}C{3.20cm}|c}
\toprule
$X$: Experiencia \textbackslash{} $Y$: Remuneración & Hasta Bs 5.000 & Más de Bs 5.000 & Suma fila \\
\midrule
Junior & $\dfrac{(30-16.50)^2}{16.50}$\newline $= 11.0455$ & $\dfrac{(23-36.50)^2}{36.50}$\newline $= 4.9932$ & $16.0386$ \\[6pt]
Consolidado & $\dfrac{(3-16.50)^2}{16.50}$\newline $= 11.0455$ & $\dfrac{(50-36.50)^2}{36.50}$\newline $= 4.9932$ & $16.0386$ \\[6pt]
\bottomrule
\end{tabular}
\end{mitabla}

\[
X_0^2 = 16.0386 + 16.0386 = 32.0772
\]

\textit{Paso 5. Estadístico teórico.} Con $f = 2$ filas y $c = 2$ columnas: $gl = (2-1)(2-1) = 1$. De la tabla chi-cuadrada, con $\alpha = 0.05$: $X^2_{1,\,0.05} = 3.841$.

\textit{Paso 6. Relación entre estadísticos y conclusión.}

\begin{center}
\begin{tabular}{l}
$32.0772 \ge 3.841 \rightarrow$ falso\\
Si ($X_0^2 < X^2_{(f-1)(c-1),\,\alpha}$)\\
\quad ``Independencia''\\
Sino\\
\quad \fbox{NO ``Independencia''}\\
FinSi
\end{tabular}
\end{center}

Se rechaza $H_0$ y se acepta $H_1$: «Remuneración» depende de «Experiencia». Porcentaje de cada fila en la columna ``Hasta Bs 5.000'': Junior $30/53 = 56.6\%$; Consolidado $3/53 = 5.7\%$.

\textit{Recomendación.} La remuneración crece con la experiencia: ofrecer a los titulados junior programas de especialización inicial que aceleren su desarrollo profesional.

